%% DoomTex :: engine/render.tex
%% Column rasteriser.
%%
%% The frame is one \hbox of 320 column \vboxes.  Each column is a stack of
%% coloured rules with \offinterlineskip, so the boxes abut exactly and the
%% vertical layout falls out of TeX's box model -- no coordinate arithmetic
%% and, crucially, no seams (verified pixel-by-pixel in spikes/s5seam.tex).
%%
%% Every band height is an integer number of DoomTex pixels, and the three
%% regions of a column are computed as differences from a running cursor, so
%% ceiling + wall + floor always sums to exactly the screen height.
\ExplSyntaxOn

\int_new:N \l_dtex_lh_int    \int_new:N \l_dtex_top_int
\int_new:N \l_dtex_ceil_int  \int_new:N \l_dtex_bot_int
\int_new:N \l_dtex_cur_int   \int_new:N \l_dtex_lvl_int
\int_new:N \l_dtex_nseg_int  \int_new:N \l_dtex_segy_int
\int_new:N \l_dtex_a_int     \int_new:N \l_dtex_b_int
\int_new:N \l_dtex_tex_int   \int_new:N \l_dtex_ucol_int

%% One horizontal band, 1 pixel wide and #3 pixels tall.
%% #4 is how far this band bleeds into its neighbour below.  The bleed is
%% taken straight back with a negative kern, so it changes what is painted
%% without moving anything.  Opaque world geometry always uses the full
%% overlap; sprites pass 0pt where the next texture row is transparent, or
%% the bleed would fill the hole the transparency is supposed to leave.
\cs_new_protected:Npn \dtexBandOv #1#2#3#4 % palette, light level, height, bleed
  {
    \int_compare:nNnT {#3} > {0}
      {
        \hbox { \textcolor { \dtexPC {#1} {#2} }
          { \vrule width \dim_eval:n { \c_dtex_px_dim + \c_dtex_overlap_dim }
                   height \dim_eval:n { \int_eval:n {#3} \c_dtex_px_dim }
                   depth #4 } }
        \kern -#4
      }
  }

\cs_new_protected:Npn \dtexBand #1#2#3
  { \dtexBandOv {#1} {#2} {#3} { \c_dtex_overlap_dim } }

%% Emit the ceiling gradient over screen rows [#1,#2), top to bottom.
\cs_new_protected:Npn \dtexCeilBands #1#2
  {
    \int_step_inline:nnn {0} {7}
      {
        \int_set:Nn \l_dtex_a_int { \dtexMax {#1} { \dtexCeilLo{##1} } }
        \int_set:Nn \l_dtex_b_int { \dtexMin {#2} { \dtexCeilHi{##1} } }
        \dtexBand { \c_dtex_ceil_pal_int } {##1}
          { \l_dtex_b_int - \l_dtex_a_int }
      }
  }

%% Emit the floor gradient over rows [#1,#2).  Floor band 7 sits at the
%% horizon and band 0 at the bottom of the screen, so this counts down.
\cs_new_protected:Npn \dtexFloorBands #1#2
  {
    \int_step_inline:nnnn {7} {-1} {0}
      {
        \int_set:Nn \l_dtex_a_int { \dtexMax {#1} { \dtexFloorLo{##1} } }
        \int_set:Nn \l_dtex_b_int { \dtexMin {#2} { \dtexFloorHi{##1} } }
        \dtexBand { \c_dtex_floor_pal_int } {##1}
          { \l_dtex_b_int - \l_dtex_a_int }
      }
  }

%% Render one screen column.
\cs_new_protected:Npn \dtexRenderColumn #1
  {
    % projected wall height for this column's distance
    \int_set:Nn \l_dtex_lh_int
      { \c_dtex_screen_h_int * 1024 / \dtexZ{#1} }
    \int_set:Nn \l_dtex_top_int
      { ( \c_dtex_screen_h_int - \l_dtex_lh_int ) / 2 }
    \int_set:Nn \l_dtex_ceil_int
      { \dtexClamp { \l_dtex_top_int } {0} { \c_dtex_screen_h_int } }
    \int_set:Nn \l_dtex_bot_int
      { \dtexClamp { \l_dtex_top_int + \l_dtex_lh_int } {0} { \c_dtex_screen_h_int } }

    % distance shading; faces perpendicular to the view go one step darker
    % so that corners read as corners even on a flat-coloured wall
    \int_set:Nn \l_dtex_lvl_int
      { \dtexClamp { ( \dtexZ{#1} - 512 ) / 3072 + \dtexSide{#1} } {0} {7} }

    \int_set:Nn \l_dtex_tex_int
      { \int_compare:nNnTF { \dtexWall{#1} } > {0} { \dtexWall{#1} } {1} }
    \int_set:Nn \l_dtex_ucol_int { \dtexU{#1} }

    % how finely to slice the wall: tall walls get the full 32 texture rows,
    % distant slivers get far fewer, which is where most of the speed is
    \int_set:Nn \l_dtex_nseg_int
      { \dtexClamp { \l_dtex_lh_int / 3 } {2} {32} }

    \vbox
      {
        \offinterlineskip
        \dtexCeilBands {0} { \l_dtex_ceil_int }
        \int_set:Nn \l_dtex_cur_int { \l_dtex_ceil_int }
        \int_step_inline:nnn {1} { \l_dtex_nseg_int }
          {
            \int_set:Nn \l_dtex_segy_int
              { \dtexClamp
                  { \l_dtex_top_int + (##1) * \l_dtex_lh_int / \l_dtex_nseg_int }
                  {0} { \c_dtex_screen_h_int } }
            \dtexBand
              { \dtexTexel { \l_dtex_tex_int } { \l_dtex_ucol_int }
                  { \dtexFloorDiv { ((##1)-1) * 32 } { \l_dtex_nseg_int } } }
              { \l_dtex_lvl_int }
              { \l_dtex_segy_int - \l_dtex_cur_int }
            \int_compare:nNnT { \l_dtex_segy_int } > { \l_dtex_cur_int }
              { \int_set:Nn \l_dtex_cur_int { \l_dtex_segy_int } }
          }
        \dtexFloorBands { \l_dtex_cur_int } { \c_dtex_screen_h_int }
      }
  }

%% The whole viewport.
\cs_new_protected:Npn \dtexRenderFrame
  {
    \hbox
      {
        \int_step_inline:nnn {0} { \c_dtex_screen_w_int - 1 }
          { \dtexRenderColumn {##1} \kern -\c_dtex_overlap_dim }
      }
  }

\ExplSyntaxOff
\ExplSyntaxOn
%% ==================================================================
%% Sprite overlay.
%%
%% Billboards are drawn in a second pass, on top of the finished frame,
%% rather than being woven into the column stacks.  Each sprite column is
%% tested against the wall distance recorded for that screen column, so
%% actors are correctly hidden behind geometry.
%%
%% The frame box has its reference point at the bottom-left, so a strip
%% whose top edge is screen row y and whose height is h is raised by
%% (screen height - y - h) pixels.
%% ==================================================================

\int_new:N \l_dtex_sprx_int  \int_new:N \l_dtex_spry_int
\int_new:N \l_dtex_sptx_int  \int_new:N \l_dtex_spty_int
\int_new:N \l_dtex_spsz_int  \int_new:N \l_dtex_spxlo_int
\int_new:N \l_dtex_spxhi_int  \int_new:N \l_dtex_spytop_int
\int_new:N \l_dtex_spid_int  \int_new:N \l_dtex_spcol_int
\int_new:N \l_dtex_sprow_int \int_new:N \l_dtex_splvl_int
\int_new:N \l_dtex_spa_int   \int_new:N \l_dtex_spb_int
\int_new:N \l_dtex_sprun_int \int_new:N \l_dtex_sppal_int

%% Inverse camera determinant.  plane is dir rotated by 90 degrees and
%% scaled by the field of view, so the determinant is exactly -fov and the
%% inverse is a constant.
\int_const:Nn \c_dtex_invdet_int { -1551 }   % 1024*1024 / -676

%% Place one 1px-wide strip of a sprite at screen column #1.  The sprite's
%% nominal top is row #2 and it is scaled to #5 pixels tall, but the strip
%% is clipped to the viewport: TeX does not clip boxes, so a sprite drawn
%% at point-blank range would otherwise spill over the HUD below.
\cs_new_protected:Npn \dtexSpriteStrip #1#2#3#4#5#6 % col, top, texcol, id, size, level
  {
    \int_set:Nn \l_dtex_spa_int { \dtexClamp {#2} {0} { \c_dtex_screen_h_int } }
    \int_set:Nn \l_dtex_spb_int
      { \dtexClamp { (#2) + (#5) } {0} { \c_dtex_screen_h_int } }
    \int_compare:nNnT { \l_dtex_spb_int } > { \l_dtex_spa_int }
      {
        \rlap
          {
            \kern \dim_eval:n { \int_eval:n {#1} \c_dtex_px_dim }
            \raisebox
              { \dim_eval:n
                  { \int_eval:n { \c_dtex_screen_h_int - \l_dtex_spb_int }
                    \c_dtex_px_dim } }
              {
                \vbox { \offinterlineskip
                  \int_step_inline:nnn {0} { \c_dtex_sprite_dim_int - 1 }
                    {
                      % rows of the strip covered by this texture row,
                      % clipped to the visible part of the sprite
                      \int_set:Nn \l_dtex_sprun_int
                        { \dtexMax { \l_dtex_spa_int }
                            { (#2) + (##1) * (#5) / \c_dtex_sprite_dim_int } }
                      \int_set:Nn \l_dtex_spcol_int
                        { \dtexMin { \l_dtex_spb_int }
                            { (#2) + ((##1)+1) * (#5) / \c_dtex_sprite_dim_int } }
                      \int_compare:nNnT { \l_dtex_spcol_int } > { \l_dtex_sprun_int }
                        {
                          \int_set:Nn \l_dtex_sppal_int { \dtexSpTexel {#4} {#3} {##1} }
                          \int_compare:nNnTF { \l_dtex_sppal_int } = {255}
                            % transparent: leave a gap so the frame shows through
                            { \kern \dim_eval:n
                                { \int_eval:n
                                    { \l_dtex_spcol_int - \l_dtex_sprun_int }
                                  \c_dtex_px_dim } }
                            {
                              \dtexBandOv { \l_dtex_sppal_int } {#6}
                                { \l_dtex_spcol_int - \l_dtex_sprun_int }
                                {
                                  \int_compare:nNnTF
                                    { \dtexSpTexel {#4} {#3} { (##1)+1 } } = {255}
                                    { 0pt } { \c_dtex_overlap_dim }
                                }
                            }
                        }
                    }
                }
              }
          }
      }
  }

%% Blit a whole sprite at a fixed screen position, with no depth test.
%% Used for the weapon and the muzzle flash, which are always in hand.
\cs_new_protected:Npn \dtexBlitSprite #1#2#3#4#5 % id, x0, y0, size, level
  {
    \int_step_inline:nnn
      { \dtexMax {#2} {0} }
      { \dtexMin { (#2) + (#4) - 1 } { \c_dtex_screen_w_int - 1 } }
      {
        \dtexSpriteStrip {##1} {#3}
          { \dtexFloorDiv { ((##1) - (#2)) * \c_dtex_sprite_dim_int } {#4} }
          {#1} {#4} {#5}
      }
  }

%% Project and draw one actor.
\cs_new_protected:Npn \dtexDrawSprite #1
  {
    \int_set:Nn \l_dtex_relx_int { \dtexEV{#1}{x} - \dtexV{px} }
    \int_set:Nn \l_dtex_rely_int { \dtexEV{#1}{y} - \dtexV{py} }
    % rotate into camera space
    \int_set:Nn \l_dtex_spa_int
      { ( \l_dtex_diry_int * \l_dtex_relx_int
          - \l_dtex_dirx_int * \l_dtex_rely_int ) / 1024 }
    \int_set:Nn \l_dtex_spb_int
      { ( 0 - \l_dtex_ply_int * \l_dtex_relx_int
          + \l_dtex_plx_int * \l_dtex_rely_int ) / 1024 }
    \int_set:Nn \l_dtex_sptx_int { \c_dtex_invdet_int * \l_dtex_spa_int / 1024 }
    \int_set:Nn \l_dtex_spty_int { \c_dtex_invdet_int * \l_dtex_spb_int / 1024 }

    % anything at or behind the camera plane, or absurdly close, is skipped
    \int_compare:nNnT { \l_dtex_spty_int } > {200}
      {
        \int_set:Nn \l_dtex_sprx_int
          { \c_dtex_screen_w_int / 2
            + ( \c_dtex_screen_w_int / 2 ) * \l_dtex_sptx_int / \l_dtex_spty_int }
        \int_set:Nn \l_dtex_spsz_int
          { \c_dtex_screen_h_int * 1024 / \l_dtex_spty_int }
        \int_compare:nNnT { \l_dtex_spsz_int } > {1}
          {
            \int_set:Nn \l_dtex_spxlo_int
              { \l_dtex_sprx_int - \l_dtex_spsz_int / 2 }
            \int_set:Nn \l_dtex_spxhi_int
              { \l_dtex_spxlo_int + \l_dtex_spsz_int }
            % actors stand on the floor, so drop them by a quarter height
            \int_set:Nn \l_dtex_spytop_int
              { ( \c_dtex_screen_h_int - \l_dtex_spsz_int ) / 2
                + \l_dtex_spsz_int / 6 }
            \int_set:Nn \l_dtex_splvl_int
              { \dtexClamp { ( \l_dtex_spty_int - 512 ) / 3072 } {0} {7} }
            \int_set:Nn \l_dtex_spid_int
              {
                \int_case:nnF { \dtexEV{#1}{type} }
                  { {1} { \int_compare:nNnTF { \dtexEV{#1}{st} } > {3} {2} {1} }
                    {2} {3} {3} {4} } {1}
              }
            \int_step_inline:nnn
              { \dtexMax { \l_dtex_spxlo_int } {0} }
              { \dtexMin { \l_dtex_spxhi_int - 1 } { \c_dtex_screen_w_int - 1 } }
              {
                % z-test against the wall recorded for this column
                \int_compare:nNnT { \l_dtex_spty_int } < { \dtexZ{##1} }
                  {
                    \dtexSpriteStrip {##1} { \l_dtex_spytop_int }
                      { \dtexFloorDiv
                          { ( (##1) - \l_dtex_spxlo_int ) * \c_dtex_sprite_dim_int }
                          { \l_dtex_spsz_int } }
                      { \l_dtex_spid_int } { \l_dtex_spsz_int } { \l_dtex_splvl_int }
                  }
              }
          }
      }
  }

%% Actors, far ones first so that nearer ones overdraw them.  With only a
%% handful of actors a selection pass is cheaper than any real sort, and it
%% keeps the ordering obviously correct.
\int_new:N \l_dtex_spsel_int  \int_new:N \l_dtex_spseld_int

%% Is this actor drawn at all?  A dead imp stays on the floor as a corpse;
%% a collected pickup is genuinely gone.
\cs_new:Npn \dtexSpriteVisible #1
  {
    \int_compare:nNnTF { \dtexEV{#1}{type} } = {1}
      { 1 }
      { \int_compare:nNnTF { \dtexEV{#1}{st} } = {5} {0} {1} }
  }

\cs_new_protected:Npn \dtex_sprite_depth:n #1
  {
    \int_set:Nn \l_dtex_relx_int { \dtexEV{#1}{x} - \dtexV{px} }
    \int_set:Nn \l_dtex_rely_int { \dtexEV{#1}{y} - \dtexV{py} }
    \int_set:Nn \l_dtex_spb_int
      { ( 0 - \l_dtex_ply_int * \l_dtex_relx_int
          + \l_dtex_plx_int * \l_dtex_rely_int ) / 1024 }
    \cs_gset:cpx { dtex_spdepth_ \int_eval:n {#1} : }
      { \int_eval:n { \c_dtex_invdet_int * \l_dtex_spb_int / 1024 } }
    \cs_gset:cpn { dtex_spdrawn_ \int_eval:n {#1} : } {0}
  }

\cs_new_protected:Npn \dtexDrawSprites
  {
    \int_step_inline:nnn {1} { \g_dtex_nent_int } { \dtex_sprite_depth:n {##1} }
    \int_step_inline:nnn {1} { \g_dtex_nent_int }
      {
        \int_set:Nn \l_dtex_spsel_int {0}
        \int_set:Nn \l_dtex_spseld_int {0}
        \int_step_inline:nnn {1} { \g_dtex_nent_int }
          {
            \bool_lazy_all:nT
              {
                { \int_compare_p:n { \use:c { dtex_spdrawn_ ####1 : } = 0 } }
                { \int_compare_p:n { \dtexSpriteVisible{####1} = 1 } }
                { \int_compare_p:n
                    { \use:c { dtex_spdepth_ ####1 : } > \l_dtex_spseld_int } }
              }
              {
                \int_set:Nn \l_dtex_spsel_int {####1}
                \int_set:Nn \l_dtex_spseld_int { \use:c { dtex_spdepth_ ####1 : } }
              }
          }
        \int_compare:nNnT { \l_dtex_spsel_int } > {0}
          {
            \cs_gset:cpn { dtex_spdrawn_ \int_use:N \l_dtex_spsel_int : } {1}
            \dtexDrawSprite { \l_dtex_spsel_int }
          }
      }
  }

\ExplSyntaxOff
